\clearpage
\fancyhead[L]{\small INSTRUMENTE.}
\SourceLeaf{415}
\begin{center}
\Large INSTRUMENTE.
\end{center}

\begin{center}
Göttingische gelehrte Anzeigen. Stück 127. Seite 1257..1267. 8. August 1818.
\end{center}

\par\indent Unsere \emph{neue Sternwarte}, dieses große und beredte Denkmal der Liebe unserer Regierung zur Wissenschaft, hat bisher in unserer Zeitschrift noch keine besondere Erwähnung gefunden, obwohl das Gebäude schon seit anderthalb Jahren so weit fertiggestellt ist, dass der Instrumentenbestand der alten Sternwarte darin untergebracht werden konnte und in einem Teil des Gebäudes Beobachtungen mit Unterbrechungen stattfinden konnten. Der Sache nach konnten es aber nur Beobachtungen von der Art sein, wie sie schon auf der alten Sternwarte möglich waren, und es schien uns nicht angebracht, darüber in dieser Zeitschrift eigens zu berichten. Mehrere davon sind bereits in astronomischen Zeitschriften veröffentlicht worden. Die neue Sternwarte, die keiner anderen nachstehen soll, kann erst mit ihren \emph{festen Meridianinstrumenten} den ihr gebührenden Rang einnehmen, und die Zeit ist nun nahe, in der sie vollständig ausgerüstet sein wird. --- Über das Äußere des Gebäudes, das der Würde seiner Bestimmung entspricht und der geschickten und geschmackvollen Arbeit unseres Universitätsbaumeisters \textsc{Müller} zur Ehre gereicht, werden wir hier nichts sagen. Auch eine vollständige Beschreibung der inneren Einrichtung bleibt einem anderen Ort vorbehalten. Aber die neuen Hauptinstrumente, die dem Gebäude erst seinen wahren wissenschaftlichen Wert geben, sollen, sobald sie nach und nach eintreffen und aufgestellt werden, in dieser Zeitschrift näher beschrieben und den Freunden der Himmelskunde bekannt gemacht werden. --- Das erste der neuen Meridianinstrumente, der \emph{Repsoldsche Meridiankreis}, kam im April dieses Jahres an und wurde vom Instrumentenbauer selbst aufgestellt. Herr \textsc{Repsold} in Hamburg hatte das Instrument bereits vor längerer Zeit ursprünglich für seinen eigenen Gebrauch angefertigt und in seiner Privatsternwarte auf dem Hamburger Wall aufgestellt. Den Freunden der Astronomie ist es nicht unbekannt, denn in der \textit{Monatlichen Correspondenz} sind mehrere Auszüge aus den Tagebüchern von Herrn \textsc{Repsold} abgedruckt. Auch in unserer Zeitschrift (1811, S. 1290 [S. 327 dieses Bandes]) haben wir über einige Beobachtungen berichtet, die Herr Professor \textsc{Schumacher} an diesem Kreis gemacht hatte. Da die Sternwarte von Herrn \textsc{Repsold} in der für Hamburg so unglücklichen Zeit abgerissen worden war und der Eigentümer das Instrument daher ungenutzt lassen musste, genehmigte unsere Regierung den Vorschlag von Herrn Hofrat \textsc{Gauss}, es für unsere neue Sternwarte anzukaufen. Der Instrumentenbauer übernahm dabei mehrere so wichtige und bedeutende Verbesserungen am Instrument, dass daraus gewissermaßen ein ganz neues geworden ist. Die Arbeit an diesen neuen Einrichtungen wurde durch manche unvorhergesehenen Umstände oft unterbrochen. Deshalb hat sich die Lieferung des Instruments so lange verzögert, obwohl auf der Sternwarte seit anderthalb Jahren alles für seine Aufstellung vorbereitet war. --- Die Bauart dieses Instruments verdient umso mehr eine ausführlichere Beschreibung, da sie ganz aus den eigenen Ideen des genialen Instrumentenbauers hervorgegangen und bis jetzt einzigartig ist. Das Instrument verbindet ein erstklassiges Durchgangsfernrohr mit einem Kreis, der dem neuen Greenwicher Mauerkreis sehr ähnlich ist. Das Fernrohr ist $7\frac{1}{2}$ Pariser Fuß lang, hat $6\frac{3}{4}$ Fuß Brennweite und 46 Linien Öffnung. Es sitzt auf einer 4 Fuß langen Achse, die wie bei anderen Durchgangsfernrohren zwischen zwei steinernen Pfeilern von $6\frac{1}{2}$ Fuß Höhe und 22 Zoll im Quadrat aufgehängt ist. Drei gesonderte Ausgleichsvorrichtungen heben sowohl die Durchbiegung des Fernrohrs als auch die der Achse auf und tragen so viel vom Gewicht des Ganzen, dass die Lagerpfannen nur mit einigen Lot belastet werden. Die Lagerpfannen \SourceLeaf{416} selbst bestehen aus Bergkristall, die vollkommen zylindrischen und gleich dicken Zapfen aus Glockenmetall. Sollten die Zapfen nach langem Gebrauch trotz der geringen Reibung etwas abgenutzt sein, lässt sich durch eine besondere Vorrichtung eine andere Stelle ihrer Oberfläche zur Auflage bringen. Das Okularrohr lässt sich parallel zur Achse verschieben, sodass jeder Faden in die Mitte des Gesichtsfeldes gebracht werden kann. Die Fäden werden durch die Achse beleuchtet. Die stärkste üblicherweise verwendete Vergrößerung ist die 96-fache. Bei günstiger Luft lassen sich Sterne bis zur dritten Größe tagsüber mühelos beobachten. Die Hängelibelle, mit der die Achse waagerecht eingestellt wird, ist äußerst sorgfältig gearbeitet und lässt Bruchteile von Sekunden sicher erkennen. Eine einzige große Platte dient beiden Pfeilern als gemeinsame Unterlage und ist ihrerseits mit großen Quadern 12 Fuß tief im Boden gegründet. Die Stabilität dieser Anordnung hat sich seit der Aufstellung aufs Beste bewährt. Azimut, waagerechte Lage der Achse und Kollimationslinie haben sich zumindest bis jetzt fast gar nicht verändert. --- Der Kreis sitzt nahe dem einen Ende fest auf der Achse und misst an der Teilung im Durchmesser $3\frac{1}{2}$ Fuß. Die Teilung wurde mit bewundernswerter Genauigkeit durch sehr saubere Striche im Abstand von jeweils 5 Minuten angebracht, während der Kreis auf der Achse selbst saß. Dadurch wurde jede Exzentrizität vermieden. Für den Versand musste der Kreis freilich wieder abgenommen werden. Beim erneuten Aufsetzen vor Ort ist nur eine fast unmerkliche Exzentrizität entstanden, die nach der sorgfältigen Prüfung von Herrn Hofrat \textsc{Gauss} $0''8$ beträgt. Abgelesen wird mit drei vortrefflichen Mikrometermikroskopen. Ihre Träger sind so an den verschiebbaren Teilen der Lager befestigt, dass ihr Abstand von der Kreisfläche vollkommen gleich bleibt, selbst wenn diese etwas von einer Ebene abweichen sollte. Ein Mikroskop sitzt unten, die beiden anderen sitzen seitlich, jeweils $90^\circ$ von ihm entfernt. Die Bauart des Instruments ließ es nicht zu, ein viertes festes Mikroskop anzubringen (früher war nur eines vorhanden). Da die beiden seitlichen Mikroskope die Wirkung einer Exzentrizität vollständig aufheben, bleibt beim Ablesen aller drei nur ein Drittel davon übrig. Es beträgt nur $0''3$ und kann berücksichtigt werden, wo es nötig erscheint. Ein viertes, \emph{bewegliches} Mikroskop zur Prüfung der Teilung selbst wird der Instrumentenbauer auf Wunsch von Herrn Hofrat \textsc{Gauss} noch nachliefern. Die Schlitten in den Mikroskopen werden durch äußerst gleichmäßig laufende Schrauben bewegt, deren Spiel durch die bekannte Erfindung von \textsc{Ramsden} beseitigt ist. Sie tragen eine kreisrunde Öffnung, in der sich die Halbierung der Teilstriche so genau erkennen lässt, dass Bruchteile von Sekunden sicher abzulesen sind. Der Kreis hat neun Speichen. Zwischen einem Paar sitzt ein Kupferzylinder, der mit einer ebenfalls ganz vortrefflichen Libelle bis auf eine halbe Sekunde genau waagerecht eingestellt werden kann. Dadurch werden etwaige Veränderungen in der Stellung der Mikroskope erkennbar. Außerdem erhält man absolute Zenitdistanzen, da sich die Neigung der Visierlinie gegen diesen Zylinder aus Beobachtungen bei umgelegtem Instrument ergibt. --- Als das Instrument ankam, war im Fernrohr noch das Fadennetz aus der ersten Herstellung eingesetzt. Es bestand aus drei senkrechten und zwei waagerechten Spinnenfäden, wobei die waagerechten $21''$ voneinander entfernt waren. Als Herr \textsc{Repsold} Ende Mai nach Göttingen zurückkehrte, ersetzte er das Fadennetz durch ein neues. Dadurch hat das Einstellen erheblich an Genauigkeit gewonnen. Das neue Fadennetz besteht aus fünf senkrechten und zwei waagerechten ausgewählten Spinnenfäden. Die waagerechten sind vollkommen parallel und stehen $12''7$ voneinander ab. Beim Beobachten bringt man den Stern in der Regel in die Mitte zwischen diese beiden Fäden. Das ist sehr genau möglich, und man kann dabei die Durchgänge an den senkrechten Fäden stets leicht und deutlich beobachten. Bei Sternen am Äquator betragen die Zeitabstände zwischen den einzelnen senkrechten Fäden etwas mehr als 16 Sekunden. --- Aus bekannten Gründen gehören die Beobachtungen des Polarsterns an jeder gut ausgestatteten \SourceLeaf{417} Sternwarte zu den täglichen Beobachtungen. Bei einem neuen Durchgangsfernrohr sind sie doppelt wichtig, weil die wichtigsten Berichtigungen auf ihnen beruhen müssen. Herr Hofrat \textsc{Gauss} hat daher alle oberen und unteren Kulminationen dieses Sterns sorgfältig beobachtet, sooft es das sehr günstige Juniwetter erlaubte. Als Beispiel dafür, was das Instrument leistet, geben wir die bisher beobachteten Rektaszensionen samt dem Vergleich mit der Hilfstafel, die Herr \textsc{Bessel} im astronomischen Jahrbuch für 1817 mitgeteilt und im folgenden Jahrgang verbessert hat. Der Bruchteil des Tages wird von der unteren Kulmination an gezählt, die zum angegebenen Datum gehört. Die letzte Spalte gibt das Quadrat der Genauigkeit jedes einzelnen Ergebnisses nach der Zahl der dazu beitragenden Fadendurchgänge an. Einige Astronomen haben dieses Quadrat sehr treffend das Gewicht der Beobachtung genannt. Als Einheit dient die Genauigkeit einer Bestimmung aus zwei aufeinanderfolgenden Kulminationen, die jeweils nur an einem Faden beobachtet wurden. Die folgenden 19 Bestimmungen beruhen auf 41 beobachteten Kulminationen, von denen jeweils zwei oder drei unmittelbar aufeinanderfolgende zu einem Ergebnis zusammengefasst wurden.

\begin{center}
\begin{tabular}{r@{\hspace{4pt}}l|c c c|c|c}
\multicolumn{2}{c}{1818} & \multicolumn{3}{c}{\shortstack{Beobachtete\\Rektaszension}} & \shortstack{Korrektur\\der Tafel} & Gewicht \\
\hline
Juni & 3.25  & $0^{\mathrm{u}}$ & $55^{\mathrm{m}}$ & $58^{\mathrm{s}}65$ & $+3^{\mathrm{s}}93$ & 2.67 \\
     & 4.25  & & & 58.19 & $+2.80$ & 5.00 \\
     & 5.25  & & & 58.56 & $+2.45$ & 5.00 \\
     & 6.25  & & & 58.91 & $+2.22$ & 4.00 \\
     & 8.50  & $0$ & $56$ & $1.14$ & $+2.97$ & 6.40 \\
     & 10.25 & & & 2.58 & $+3.20$ & 3.75 \\
     & 11.25 & & & 2.68 & $+2.59$ & 1.67 \\
     & 12.50 & & & 3.31 & $+2.36$ & 4.29 \\
     & 16.75 & & & 6.11 & $+2.13$ & 4.44 \\
     & 18.25 & & & 7.31 & $+2.26$ & 5.00 \\
     & 19.25 & & & 8.29 & $+2.51$ & 2.86 \\
     & 21.25 & & & 10.49 & $+3.26$ & 5.00 \\
     & 22.75 & & & 12.02 & $+3.60$ & 1.33 \\
     & 26.25 & & & 13.49 & $+2\ 60$ & 1.60 \\ 
     & 27.25 & & & 14.77 & $+3.14$ & 4.44 \\
     & 29.50 & & & 16.71 & $+3.42$ & 3.43 \\
Juli & 7.50  & & & 22.51 & $+3.31$ & 6.00 \\
     & 8.25  & & & 22.97 & $+3.22$ & 1.67 \\
     & 9.25  & & & 23.89 & $+3.40$ & 3.75 \\
\end{tabular}
\end{center}

\par\indent Der Mittelwert der aus diesen 19 Beobachtungen ermittelten Korrekturen der Tafel ist $+2^{\mathrm{s}}88$. Wenn man sich gestattet, den Tafelfehler während dieser Zeit als konstant anzusehen, ergibt die Anwendung der Wahrscheinlichkeitstheorie von Herrn Hofrat \textsc{Gauss} für den wahrscheinlichen Fehler einer Beobachtung, deren Genauigkeit 1 ist, den Wert $=0^{\mathrm{s}}64$ und für den wahrscheinlichen Fehler des Endergebnisses den Wert $=0^{\mathrm{s}}076$. Als zweites Beispiel fügen wir die Beobachtungen des Uranus bei, die Herr Hofrat \textsc{Gauss} zur Zeit der Opposition dieses Planeten angestellt hat. Die Deklinationen beruhen auf dem Vergleich mit der unmittelbar vorausgegangenen unteren und der unmittelbar folgenden oberen Kulmination des Polarsterns. Ausgenommen sind nur die Beobachtungen vom 7. und 13. Juni, an denen nur die untere Kulmination beobachtet worden war und die Deklination des Polarsterns aus \textsc{Bessels} Tafel zugrunde gelegt wurde. Die Refraktionen sind \textsc{Bessels} Tafel entnommen, die auf den Beobachtungen von \textsc{Bradley} beruht. Würde man sie \textsc{Carlinis} Tafel entnehmen, so würden die Deklinationen des Uranus $3''$ weiter südlich ausfallen. Sobald in Zukunft genügend geeignete Beobachtungen vorliegen, wird sich zeigen, welche der beiden Tafeln dem hiesigen Klima am besten entspricht.

\SourceLeaf{418}
\begin{center}
\setlength{\tabcolsep}{3pt}
\begin{tabular}{l|ccc|ccc|ccc}
\hline
1818 & \multicolumn{3}{c|}{Mittlere Zeit} & \multicolumn{3}{c|}{Rektaszension} & \multicolumn{3}{c}{Südl. Dekl.} \\
\hline
Juni 3 & $12^{\mathrm u}$ & $22^{\mathrm m}$ & $17^{\mathrm s}5$ & $257^\circ$ & $23'$ & $28''3$ & $23^\circ$ & $1'$ & $41''5$ \\
4 & 12 & 18 & 10.9 & & 20 & 47.8 & & 1 & 29.0 \\
5 & 12 & 14 & 4.6 & & 18 & 9.9 & & 1 & 14.4 \\
7 & 12 & 5 & 51.6 & & 12 & 54.4 & & 0 & 50.5 \\
8 & 12 & 1 & 45.1 & & 10 & 15.0 & & 0 & 38.4 \\
9 & 11 & 57 & 38.6 & & 7 & 36.4 & & 0 & 27.4 \\
10 & 11 & 53 & 32.1 & & 4 & 56.5 & & 0 & 12.6 \\
11 & 11 & 49 & 25.7 & & 2 & 18.0 & & 0 & 2.9 \\
12 & 11 & 45 & 19.2 & 256 & 59 & 33.9 & 22 & 59 & $50\;9$ \\ 
13 & 11 & 41 & 12.7 & & 57 & 0.4 & & 59 & 40.5 \\
\hline
\end{tabular}
\end{center}

\Needspace{18\baselineskip}\par\indent Herr Hofrat \textsc{Gauss} übertrug den Vergleich dieser Beobachtungen mit den Uranustafeln von \textsc{Delambre} Herrn \textsc{Dirksen}, der sich bei uns mit besonderem Eifer dem Studium der Astronomie widmet. Die Ergebnisse dieses Vergleichs sind folgende:

\begin{center}
\begin{tabular}{c|cc}
& \multicolumn{2}{c}{Unterschied} \\
& Rektaszension & Deklination \\
\hline
Juni 3 & $-54.1$ & $+10.3$ \\
4 & $-52.1$ & $+10.6$ \\
5 & $-52.9$ & $+12.9$ \\
7 & $-55.8$ & $+12.3$ \\
8 & $-55.9$ & $+12.0$ \\
9 & $-56\;8$ & $+10.5$ \\
10 & $-56.4$ & $+12.9$ \\
11 & $-57.2$ & $+10.1$ \\
12 & $-52.4$ & $+9\;4$ \\ 
13 & $-58.1$ & $+7.4$ \\
\hline
Mittel & $-55.2$ & $+10.8$ \\
\end{tabular}
\end{center}

\Needspace{15\baselineskip}\par\indent Daraus ergab sich nach der Rechnung von Herrn \textsc{Dirksen}:

\begin{center}
Zeit der Opposition am 9. Juni 1818.\\
$5^{\mathrm u}\,30^{\mathrm m}\,43^{\mathrm s}$ Mittlere Zeit in Göttingen
\end{center}
\begin{center}
\begin{tabular}{ll}
Wahre Länge . . . . . . . . . . . . . . & $258^\circ\;10'\;34''5$ \\
Geozentrische Breite . . . . . . . . . . . . & $0\;4\;12.1$ \\
Heliozentrische . . . . . . . . . . . . . . & $0\;3\;58.7$ Südl. \\
\end{tabular}
\end{center}

\noindent und der Fehler der Tafeln von \textsc{Delambre}

\begin{center}
\begin{tabular}{ll}
in der Länge . . . . . . . . . . . . . . & $-45''3$ \\
in heliozentrischer Breite . . . . . . . . . . & $+14.6$ \\
\end{tabular}
\end{center}

\par\indent Wir können diesen Bericht nicht schließen, ohne auch die Pendeluhr von \textsc{Shelton} zu erwähnen. Sie war auf der alten Sternwarte bereits 47 Jahre in Gebrauch gewesen, aber erst jetzt lässt sich ihre Güte voll würdigen, da die Sternwarte früher über kein Mittel zur genauesten Zeitbestimmung verfügte. Zum Beleg fügen wir nur die Aufzeichnungen über ihren Gang während eines Monats bei. Wir bemerken dazu, dass der Gang während der gesamten Zeit, in der das \textsc{Repsold}sche Instrument zu seiner Bestimmung diente, ebenso gleichmäßig war. Das Datum bezieht sich auf Sonnentage. Zur besseren Übersicht sind die Stunden der Uhrzeit jeweils über 24 hinaus bis zum nächsten Mittag weitergezählt. Die letzte Spalte enthält die Anzahl der Sterne aus \textsc{Maskelynes} Fundamentalkatalog, auf denen die Zeitbestimmung beruht.

\SourceLeaf{419}
\begin{center}
\begin{tabular}{ll|c|c|c}
 & & Uhrzeit & \shortstack{Abweichung von\\der Sternzeit} & \shortstack{Verglichene\\Sterne} \\
\hline
Juni & 15 & $15^{\mathrm u}\;22^{\mathrm m}$ & $+0^{\mathrm s}41$ & 4 \\
 & 17 & $20\;\;50$ & $+0.28$ & 2 \\
 & 18 & $11\;\;37$ & $+0.29$ & 2 \\
 & 20 & $12\;\;28$ & $+0.11$ & 2 \\
 & 21 & $11\;\;37$ & $-0.06$ & 2 \\
 & 22 & $7\;\;\;30$ & $-0.59$ & 1 \\
 & 23 & $11\;\;40$ & $-0.90$ & 1 \\
 & 26 & $14\;\;54$ & $-1.78$ & 5 \\
 & 27 & $16\;\;15$ & $-1.86$ & 4 \\
 & 28 & $14\;\;41$ & $-1.99$ & 1 \\
 & 29 & $18\;\;\;\;1$ & $-2.32$ & 6 \\
 & 30 & $16\;\;\;\;9$ & $-2.43$ & 5 \\
\end{tabular}
\end{center}

\begin{center}
\begin{tabular}{ll|c|c|c}
 & & Uhrzeit & \shortstack{Abweichung von\\der Sternzeit} & \shortstack{Verglichene\\Sterne} \\
\hline
Juli & 1 & $14^{\mathrm u}\;7^{\mathrm m}$ & $-2^{\mathrm s}92$ & 1 \\
 & 2 & $14\;\;51$ & $-3.11$ & 2 \\
 & 6 & $25\;\;17$ & $-3.87$ & 4 \\
 & 7 & $22\;\;31$ & $-3.90$ & 7 \\
 & 8 & $16\;\;36$ & $-4.03$ & 4 \\
 & 9 & $28\;\;\;7$ & $-4.21$ & 4 \\
 & 10 & $15\;\;\;3$ & $-4.28$ & 3 \\
 & 13 & $29\;\;24$ & $-4.53$ & 2 \\
 & 14 & $29\;\;14$ & $-4.62$ & 4 \\
 & 15 & $15\;\;56$ & $-4.73$ & 5 \\
 & 16 & $14\;\;41$ & $-4.93$ & 1 \\
\end{tabular}
\end{center}

\par\indent Dabei muss auch erwähnt werden, dass die Uhr vorläufig aufgestellt ist. Sie ruht zwar auf einem eigenen steinernen Fundament, aber nur über ihr Gehäuse, und ist nicht unmittelbar an dem steinernen Pfeiler befestigt, neben dem das Gehäuse steht. Sie läuft mit einem einzigen Aufziehen einen Monat lang. Etwa in der Mitte dieser Zeit steht das Gewicht der Pendellinse gegenüber. Das hat man früher wohl für nachteilig gehalten und durch mancherlei Vorrichtungen zu vermeiden gesucht. Die Aufzeichnungen über die \textsc{Shelton}sche Uhr zeigen aber um diese Zeit keinerlei spürbare Veränderung des Gangs. Überhaupt ist ihr Gang viel regelmäßiger als der mancher anderen Uhren auf den bedeutendsten Sternwarten, auch solcher, die unmittelbar an einem steinernen Pfeiler befestigt sind. Sollte man daraus nicht schließen, dass beide Umstände von viel geringerer Bedeutung sind, als man bisher gewöhnlich geglaubt hat?
